\section{Introduction}
\label{sec:intro}
\IEEEPARstart{R}{ibs} and clavicles overlap the lung fields on a chest radiograph and can hide small lesions. Dual-energy subtraction
(DES) separates bone from soft tissue, but it needs dedicated hardware and a second exposure \cite{kuhlman2006des,sun2026bsldm}.
Software bone suppression avoids both: radiologists find more nodules with suppressed images \cite{freedman2011lung}, and suppressed
images have been compared with computer-aided detection \cite{schalekamp2014cad} and used as classifier inputs
\cite{rajaraman2021bone}. Learning-based suppressors need paired supervision. It can come from real DES pairs
\cite{chen2024bsdiff,sun2026bsldm,sun2025gllcm,sun2026debonedit}, which are scarce, or from chest CT, whose bone and soft-tissue
components can be projected separately and are therefore known exactly \cite{angaitkar2026anatomy}. Public projections of CT-RATE
volumes \cite{hamamci2026ctrate} make the second route practical.

A suppressor judged by image fidelity can nevertheless erase lesions. A network rewarded for predicting bone can move lesion signal into
the bone image, so that the lesion disappears from the suppressed radiograph, and PSNR or SSIM computed over the whole image are dominated
by normal anatomy. Because a model that does nothing never absorbs a lesion, absorption has to be read next to suppression strength.
Simulated nodules have been used to measure contrast after hand-built clavicle suppression and in phantoms
\cite{hogeweg2013suppression,takagi2020dose}, and learned image-to-image models are known to remove or invent content in other
modalities \cite{cohen2018distribution,antun2020instabilities,bhadra2021hallucinations,tivnan2024hallucination}. In our search (PubMed,
arXiv and Crossref) we found no audit of learned bone suppression against an exactly known, additive lesion signal, and no training
objective aimed at it.

Whether suppression helps prediction is a separate question. The suppressed images are computed from the radiograph and add no
information that a classifier could not extract itself, so a gain must come from the prior that the suppressor learned from CT or from the
way a network uses the extra images. A network that reads three images spends three times the computation, so a gain can be credited to
suppression only against a control of the same capacity, and seed variation and analytic flexibility can easily produce spurious
differences \cite{varoquaux2022failures,simmons2011falsepositive}. Every comparison here is therefore paired, repeated over seeds and
fixed in advance.

We ask whether CT-supervised bone suppression can preserve pulmonary lesions and improve nodule detection, focal disease classification
and localisation. Our contributions are:
\begin{itemize}
\item A counterfactual lesion audit that measures the fraction $A$ of an exactly known, additive lesion signal moved into the bone image,
 at sites shared by all models and reported beside suppression strength.
\item A twin lesion-invariance regulariser, trained on public CT projections in a three-model ensemble. On VinDr-RibCXR the ensemble
 absorbed 0.051 of an added lesion, against 0.149 for a released CT-supervised model and 0.178 for a matched control, at non-inferior
 bone-region PSNR (21.46 vs.\ 21.51\,dB). The effect replicated across three encoders, on JSRT and on paediatric radiographs
 ($A=0.084$ vs.\ 0.217 and 0.251).
\item Evidence that suppressed images help prediction when added as extra input channels: Node21 detection (CPM $+0.028$,
 95\,\% CI 0.009--0.047), VinBigData detection ($+0.008$, 0.002--0.014) and NIH macro-AUROC ($+0.006$, 0.003--0.008), whereas the
 soft-tissue image alone gave no gain.
\item A fusion ablation of five input configurations (175 runs, three data sets, a capacity-matched control). A three-stream network gave
 no overall gain; the benefit was concentrated in focal findings (NIH soft-tissue focal AUROC $+0.0089$, Holm $p=0.0024$; nodule
 $+0.0129$, Holm $p=0.00121$) and in the localisation of focal lesions (pointing-game interaction $+0.145$, $p=0.0002$), at a cost for
 diffuse findings.
\item Deployment analyses: ensemble uncertainty that needs recalibration on real radiographs, a paediatric external validation and an
 exploratory site calibration of the fidelity metric ($+2.12$\,dB in the bone region without raising lesion absorption).
\end{itemize}
Fig.~\ref{fig:workflow} summarises the two studies. Section~\ref{sec:related} reviews related work, Section~\ref{sec:methods} describes
the data, the audit, the suppressors and the downstream studies, Section~\ref{sec:results} reports the results, and
Section~\ref{sec:discussion} discusses them together with their limitations.

\begin{figure*}[!t]
\centering
\safefig[0.9\textwidth]{figures/fig_tmi_workflow.pdf}
\caption{Overview. Study~I trains CT-supervised bone suppressors with a twin lesion-invariance regulariser, audits their lesion absorption
and tests the suppressed images on detection and classification; Study~II compares five input configurations of a multi-stream
classifier, including a capacity-matched control.}
\label{fig:workflow}
\end{figure*}
